\section{Introduction}

A standard assumption in regression is that feature omission is simply due to noise. However, in many real-world regression tasks, the sender has a preference over the regression outcomes. Because features are the sender's private information that the receiver (regressor) cannot directly observe, this information asymmetry allows the sender to selectively reveal only those features that get them closer to their desired outcome. The receiver must therefore account for this strategic behavior. This task is inherently challenging because the receiver cannot initially determine whether a missing feature is absent due to random channel noise or deliberate omission. Furthermore, even if the receiver can deduce that an omission was intentional, identifying the direction of the underlying value remains difficult if the sender has non-uniform preferences over the outcome.

To illustrate this problem, consider a policy planner collecting public health data. A citizen with poor health indicators might deliberately leave a column blank to present a lower risk profile if they want to purchase health insurance in the future. Conversely, an individual with good health indicators might also withhold their data to present a high-risk profile if they want to maximize a payout claim. This difference in motivations for leaving the column blank represents the non-uniform preferences of the senders. For the policymaker, the regression task is difficult not only due to the ambiguity of the omission's cause but also due to the ambiguous direction of the underlying feature value, as both high-risk and low-risk individuals have distinct incentives to conceal their data.

Even though strategic withholding complicates the regression task, it is not always unfavourable to the receiver's performance. We show that there exist environments where voluntary feature revelation actually improves prediction accuracy over a mandatory feature-revelation baseline.  
Nonetheless, our strategic-risk upper bound has a larger complexity term than the mandatory feature-revelation baseline.

\subsection{Contributions}

In this chapter, we study the problem of strategic feature disclosure for regression. We work under the structural assumption of an all-or-nothing revelation model, where the sender either reveals the feature vector or withholds it from the receiver. We examine three problems based on the information structure available to the receiver.

First in the \textit{Decision Problem}, where the environment is known to the receiver, we define the strategic advantage $\Delta(\mathcal{E})$ as the receiver's performance gain from voluntary feature revelation over the mandatory baseline, and derive a necessary and a sufficient conditions on the environment for this advantage to be positive.

Second in the \textit{Offline Learning Problem}, where the receiver has access to a finite, fully revealed training dataset alongside the sender's preferences, we propose an ERM-based algorithm that estimates the strategic advantage and achieves a uniform generalization bound of $\cO\left(\frac{d^3 \sqrt{\log n}}{\sqrt{n}} + d^2\sqrt{\frac{\log(1/\delta)}{n}}\right)$, where $d$ is the dimension of the feature space.

Third in the \textit{Online Learning Problem}, where the receiver learns sequentially under bandit feedback, we show that the expected risk functions arising in our setting lie in a Sobolev space $H^{\nu+d/2}$ with $\nu=5/2$ whenever the joint density $p_{U,X,Y}$ is smooth and has compact support. Thus, the regularity assumptions underlying the standard Matérn-kernelized Gaussian-process Thompson sampling (GP-TS) algorithm are satisfied in our setting, allowing us to apply GP-TS over the joint action space of regimes and regression parameters and obtain its regret guarantee of $\widetilde{\cO}\bigl(T^{\frac{5+3d}{10+2d}}\bigr)$.

\subsection{Related Work}

% Our work lies in learning with strategic agents~\cite{bruckner2011stackelberg,hardt2016strategic,dong2018strategic}, typically modeled as sender--receiver games. Prior work studies robustness to manipulation~\cite{zhang2015adversarial,chen2018strategyproof}, causal shifts~\cite{shavit2020causal}, incentive design~\cite{jin2022incentive}, and strategic precision degradation~\cite{gast2020linear}. A common assumption is that senders may arbitrarily falsify features, which excludes settings where only verifiable features can be omitted. \cite{krishnaswamy2021classification} study withholding but with uniform preferences in classification; \cite{sundaram2023pac,singh2025strategic} allow non-uniform preferences yet still permit alteration. Classic disclosure models~\cite{grossman1981informational,milgrom1981good} treat voluntary revelation with the sender as leader. To our knowledge, combining strategic withholding, non-uniform preferences, and a receiver-led regression game remains open; our conference version~\cite{singh2026linear} develops the known-environment and offline-learning foundations.

This chapter contributes to the literature on learning in the presence of strategic agents \cite{bruckner2011stackelberg,hardt2016strategic,dong2018strategic}, where interactions are commonly formulated as sender--receiver games. Existing studies have investigated robustness to strategic manipulation \cite{zhang2015adversarial,chen2018strategyproof}, distributional changes induced by causal interventions \cite{shavit2020causal}, incentive mechanisms \cite{jin2022incentive}, and strategic precision degradation \cite{gast2020linear}. Most of this literature assumes that the sender can freely manipulate the disclosed information and has uniform preference. In contrast, we study environments in which disclosed features are verifiable, and the sender's strategic choice is limited to withholding information, while also allowing non-uniform preferences. Among the closest related works, \cite{krishnaswamy2021classification} considers feature withholding but assumes uniform preferences and restricts attention to classification. The works of \cite{sundaram2023pac,singh2025strategic} allow non-uniform preferences but instead permit strategic modification of features. In the economics literature, voluntary disclosure has been extensively studied \cite{grossman1981informational,milgrom1981good}; however, these models typically treat the sender as the Stackelberg leader. A preliminary version of the results in this chapter appeared in the conference paper~\cite{singh2026linear}, which developed the decision and offline-learning foundations. This chapter is derived from our manuscript~\cite{singh2026revelation_journal}, which provides complete proofs and additionally investigates the online continuous-bandit learning setting. The online learning component of this chapter is essentially an application of Gaussian-process Thompson sampling for continuous bandits \cite{srinivas2010gp,chowdhury2017kernelized,vakili2021information}. While the algorithm itself is standard, its regret guarantees require the expected risk functions to lie in an appropriate RKHS. We establish that this premise holds in our setting for environments with smooth joint densities over compact support, thereby justifying the application of the existing theory. To the best of our knowledge, the setting of jointly considering strategic feature withholding, non-uniform preferences, and a receiver-led game has not been previously studied, and was first explored in our conference work~\cite{singh2026linear}.


